\section{Introduction}
\label{sec:introduction}

Direct lithium extraction from unconventional brines requires a thermodynamic and process model that can separate the lithium recovery question from the brine-severity question. Produced waters are especially relevant because they can combine high total dissolved solids, large sodium excess, and measurable lithium concentrations with existing water-handling infrastructure. Those same features make the separation problem hard: the extraction solvent sees a high-ionic-strength feed, and any direct solvent model must either include divalent competition explicitly or define a pretreated feed boundary.

This manuscript frames the case study around a pretreated \Smackover \LiIon/\NaIon stream. The current presentation and analysis chain does not use the older HBTA/TOPO solvent-selection lane as the active case-study basis. Instead, it uses the hydrophobic \DES plus \TOPO extraction chemistry reported by Rezaee and co-workers as the solvent evidence base \cite{Rezaee2025RSM,Rezaee2026PCSASFTSI}. The resulting contribution is not a new extraction solvent, but a traceable engineering workflow that connects literature extraction data, ePC-SAFT-facing parameter artifacts, a distribution-coefficient surrogate, and downstream \Prommis/\IDAES process screening.

The first-draft paper is organized around four questions:
\begin{enumerate}
    \item Which produced-water composition should be used as the clean \LiIon/\NaIon case-study basis?
    \item How should Rezaee extraction data be converted into process-facing distribution coefficients without claiming unsupported direct reactive-LLE closure?
    \item What range of lithium recovery, sodium co-extraction, selectivity, and transfer coefficients is produced by the current calibrated screening matrix?
    \item Which assumptions must be promoted before the \Prommis/\IDAES costing layer can be interpreted as a final techno-economic analysis?
\end{enumerate}
